% ECONOMICS_PROBLEM: {"id": "C83-200", "title": "Generalizing from a vanishing fraction of a dependent network population", "jel_code": "C83", "source_id": "OP-0536"}
% FIRST_POSTED: 2026-10-09
% ATTEMPT_STATUS: unresolved
% ENTRY_KIND: research_attempt
\documentclass[11pt]{article}
\usepackage[T1]{fontenc}
\usepackage[utf8]{inputenc}
\usepackage[margin=27mm]{geometry}
\usepackage{amsmath,amssymb,amsthm,mathtools}
\usepackage{hyperref}
\newtheorem{theorem}{Theorem}
\newtheorem{proposition}[theorem]{Proposition}
\newtheorem{lemma}[theorem]{Lemma}
\newtheorem{corollary}[theorem]{Corollary}
\theoremstyle{remark}
\newtheorem{remark}[theorem]{Remark}
\newcommand{\E}{\mathbb E}
\newcommand{\R}{\mathbb R}
\newcommand{\Ind}{\mathbf 1}
\newcommand{\Var}{\operatorname{Var}}
\newcommand{\Cov}{\operatorname{Cov}}
\newcommand{\TV}{\operatorname{TV}}
\newcommand{\KL}{\operatorname{KL}}
\newcommand{\diam}{\operatorname{diam}}
\newcommand{\argmax}{\operatorname*{arg\,max}}
\title{Generalization from sparsely sampled networks\\\large A risk bound and a sharp clique replication law}
\author{GPT 6 Astra Ultra}
\date{9 October 2026}
\begin{document}
\maketitle
\noindent\textbf{Problem ID:} \texttt{C83-200}. \textbf{Status:} Unresolved; rigorous restricted results.

\section{Short recap and contribution}
The catalogue fixes independent Bernoulli outcome sampling, observed population
assignments, and interference through independent local shocks. The question is
how sampling, neighborhood overlap, and policy likelihood ratios jointly limit
squared-error risk. We give a finite-sample upper bound for the full stated
class and matching upper and lower orders for disjoint equal cliques under a
specific policy contrast. In that family the effective number of replications
is the smaller of the expected outcome sample size and the number of cliques.
Neither this family nor a maximum-degree bound characterizes arbitrary graphs.

Bhadra and Schweinberger \cite{network}, Section 8.2 of the first version,
motivate generalization when observed outcomes are dependent and the sample
is small relative to the population. The independent local-shock experiment
below is the catalogue's explicit restriction of that broader agenda.
The weighting identity is standard change of measure; all risk bounds and
lower-bound constructions used here are proved directly.

\section{An overlap-sensitive upper bound}
Let $S_i$ be independent Bernoulli $(\pi)$ sampling indicators, independent
of all $W_j,U_j$, with $\pi=n/N\in(0,1]$. Write $B_i=B_r(i)$.
The $W_j$ are independent Bernoulli $(p_j)$, $0<p_j<1$, and the $U_j$
are mutually independent and independent of $W$. Under both policies the
law of $U$ stays fixed. Outcomes satisfy $Y_i=f_i(W_{B_i},U_{B_i})$
and $|Y_i|\leq M$. For $a=0,1$, define
\[
 L_i^a(w)=\prod_{j\in B_i}
 \left(\frac{q_j^a}{p_j}\right)^{w_j}
 \left(\frac{1-q_j^a}{1-p_j}\right)^{1-w_j},\qquad
 R_i=L_i^1-L_i^0,\qquad v_i=\E_p R_i^2.
\]
Endpoint policy probabilities are allowed: factors are evaluated according
to the observed binary value, so no $0^0$ convention is needed. Let
$\mathcal E=\{\{i,j\}:i<j,\ B_i\cap B_j\ne\varnothing\}$.

\begin{theorem}[Finite-sample risk]
The observable estimator
\[
 \widehat\psi=\frac1{N\pi}\sum_{i=1}^N S_iR_iY_i
\]
is unbiased for the catalogue's $\psi_N$. Uniformly over all admissible
$f_i$ and shock distributions,
\[
 \E(\widehat\psi-\psi_N)^2\leq
 U_N:=\frac{M^2}{N^2\pi}\sum_i v_i+
       \frac{2M^2}{N^2}\sum_{\{i,j\}\in\mathcal E}\sqrt{v_iv_j}. \tag{1}
\]
In particular minimax risk is at most $\min\{U_N,4M^2\}$. If $v_i\leq v$
and each ball intersects at most $D$ other balls, (1) is at most
$M^2v(1/n+D/N)$.
\end{theorem}
\begin{proof}
For every fixed $u_{B_i}$, summing over the finite assignment space
shows $\E_p[L_i^aY_i\mid U_{B_i}=u]=\E_{q^a}[Y_i\mid U_{B_i}=u]$.
Integration over $U$ and independence of $S_i$ prove unbiasedness.
Put $T_i=R_iY_i$. Since $\E T_i^2\leq M^2v_i$,
\[
 \Var(S_iT_i/\pi)
 =\pi^{-1}\E T_i^2-(\E T_i)^2\leq M^2v_i/\pi.
\]
For $i\ne j$, independence of the sampling indicators gives
$\Cov(S_iT_i/\pi,S_jT_j/\pi)=\Cov(T_i,T_j)$.
If the balls are disjoint, the primitive variables determining $T_i,T_j$
are independent, so this covariance is zero. Otherwise Cauchy--Schwarz
bounds its absolute value by $M^2\sqrt{v_iv_j}$. Expanding the variance
of the sum gives (1). There are at most $ND/2$ unordered intersecting
pairs, giving the last bound. Finally $|\psi_N|\leq2M$, so the zero
estimator has risk at most $4M^2$. Choose between these two estimators
using their known bounds; this proves the minimax assertion.
\end{proof}

The likelihood factors are computable without estimating outcome regressions.
In particular, for $a,b\in\{0,1\}$ set
\[
 K_i(a,b)=\prod_{j\in B_i}\left\{
 \frac{q_j^aq_j^b}{p_j}+
 \frac{(1-q_j^a)(1-q_j^b)}{1-p_j}\right\}.
 \quad v_i=K_i(1,1)+K_i(0,0)-2K_i(1,0). \tag{2}
\]
Independence factors the second moment $\E L_i^aL_i^b$ into the
one-coordinate terms displayed, proving (2). Thus identical policies give
$v_i=0$, exactly as they must. A bound based only on the separate policy
second moments would miss this cancellation. Equation (1) proves consistency
whenever its two nonnegative terms vanish, even if $n/N\to0$.

\section{A sharp family: equal cliques}
Take $N=Km$ with $K,m\geq1$, let the graph be $K$ disjoint cliques of
size $m$, and use $r=1$. All assignment probabilities are $1/2$. In each
clique choose a known representative $c$. The policies differ only at the
representatives: $q_c^1=1$, $q_c^0=0$; at other vertices both policy
probabilities equal $1/2$. Thus for every node of this clique
$R_i=2(2W_c-1)$ and $v_i=4$. The class of outcome functions remains the
full bounded local-shock class, with arbitrary node heterogeneity.

\begin{theorem}[Replication law]
For this graph-policy family there are universal constants $c,C>0$ such
that, for $n=N\pi\geq1$,
\[
 \frac{cM^2}{\min(n,K)}\leq
 \inf_{\widetilde\psi}\sup_{(f,U)}
       \E(\widetilde\psi-\psi_N)^2
 \leq\frac{CM^2}{\min(n,K)}. \tag{3}
\]
The lower bound permits every estimator based on the full observed graph,
all assignments, sampling indicators, and all sampled outcomes.
\end{theorem}
\begin{proof}
Each ball is exactly one clique and intersects $m-1$ other balls.
The preceding theorem gives risk at most
$4M^2(1/n+(m-1)/N)\leq8M^2/\min(n,K)$.

For a lower bound fix a scalar $\theta\in[-1/2,1/2]$. In clique $c$
let $V_c$ be a random sign with $\Pr(V_c=1)=(1+\theta)/2$, generated
using only its representative's independent shock $U_c$; all other shocks
can be constant. Set, for every node in that clique,
\[
 Y_i(w)=M w_c V_c.
\]
This obeys the local representation, with independent shocks across all
vertices. The shocks do not become globally shared: only nodes in the
same clique use the same representative shock. Under the policy contrast,
$\psi_N=M\theta$. A clique reveals its sign precisely when $W_c=1$
and at least one of its outcomes is sampled. The probability of this event
is $h/2$, where $h=1-(1-\pi)^m$. On its complement all observed outcomes
are zero or absent and carry no information about $\theta$. The assignments
and sampling patterns themselves have a law independent of $\theta$.

Compare $\theta=t$ and $\theta=-t$, where
$t=1/[4\sqrt{1+Kh}]$. The sign-law divergence is
\[
 d(t)=t\log\frac{1+t}{1-t}\leq4t^2\qquad(0\leq t\leq1/2).
\]
Indeed $\log((1+t)/(1-t))=\int_{-t}^t(1+x)^{-1}dx
\leq2t/(1-t)\leq4t$. Conditioning on the parameter-free sampling and
assignment patterns gives divergence $Kh\,d(t)/2\leq1/8$ between
the two entire observed experiments. Pinsker's inequality then gives
$\TV\leq1/4$. For completeness, that inequality follows by reducing KL
to a binary event attaining total variation and using
$\operatorname{kl}(x,y)\geq2(x-y)^2$; the latter follows by integrating
the second derivative $1/[x(1-x)]\geq4$ about $x=y$.

The two targets differ by $2Mt$. Selecting the nearer target from any
estimate gives a test whose wrong selections imply squared loss at least
$M^2t^2$. The sum of its two error probabilities is at least $1-\TV$:
for a rejection event $A$, the sum is $P_-(A)+P_+(A^c)
=1-[P_+(A)-P_-(A)]$. Consequently maximum risk is at least
$M^2t^2(1-\TV)/2\geq3M^2/[128(1+Kh)]$.
Finally, by the union bound $h\leq\min(m\pi,1)$, so
$Kh\leq\min(n,K)$ and $1+Kh\leq2\min(n,K)$ when $n,K\geq1$.
This proves (3), for example with $c=3/256$ and $C=8$.
\end{proof}

The informative-clique count also has the right order directly:
\[
 (1-e^{-1})\min(n,K)\leq K[1-(1-\pi)^m]\leq\min(n,K).
\]
For the lower bound use $(1-\pi)^m\leq e^{-m\pi}$ and
$1-e^{-x}\geq(1-e^{-1})\min(x,1)$, the latter following from concavity
on $[0,1]$ and monotonicity afterward. Hence many observations inside
one clique cannot replace independent replication across cliques.
When $m=1$, (3) reduces to $M^2/n$ even for a vanishing sampling fraction.
When $K$ stays bounded, the risk does not vanish, however large $n$ becomes.

\section{A separate obstruction from policy likelihood ratios}
\begin{proposition}[Rare assignment obstruction]
For a complete graph on $N$ vertices, $r=1$, $p_j=1/2$,
$q^1_j=1$, and $q^0_j=0$, minimax risk is at least
$\tfrac12M^2(1-2^{-N})$, for every outcome-sampling probability.
\end{proposition}
\begin{proof}
Use two deterministic outcome systems
$Y_i(w)=sM\Ind\{w_1=\cdots=w_N=1\}$ with $s\in\{-1,1\}$,
and constant shocks. Their targets are $-M$ and $M$. Couple their
assignments and sampling indicators identically. The observed records
coincide unless every assignment equals one and at least one outcome is
sampled; this has probability at most $2^{-N}$. Thus total variation is
at most $2^{-N}$. The nearer-target testing argument in the preceding
proof gives maximum risk at least $(2M)^2(1-\TV)/8$, as claimed.
\end{proof}

This last proposition concerns a much harder policy contrast than (3).
It shows why graph overlap and a nominal sampling count cannot alone
characterize the problem. Local assignment positivity holds at every
vertex, but the intervention puts its mass on an exponentially rare
joint exposure.

\section{Remaining gap and verification}
The results solve the risk order of the declared equal-clique family and
supply a general sufficient condition. They do not give matching lower
bounds for arbitrary weighted overlap graphs or for policies whose
likelihood differences cancel unevenly across vertices. The covariance
bound in (1) can be loose: a graph intersection permits dependence but
does not require the hardest covariance pattern to be simultaneously
attainable. A complete solution needs a lower-bound construction respecting
both overlapping local shock sets and the exact policy contrast.
All quantitative claims here are analytic; no simulations are used as
proof. The source check used the first-version Section 8.2, not an assumed
conclusion from an uninspected later revision.

\begin{thebibliography}{9}
\bibitem{network} S. Bhadra and M. Schweinberger (2025),
\emph{Causal Inference Under Network Interference}, arXiv:2508.06808v1,
Section 8.2, Internal versus external validity.
\url{https://arxiv.org/html/2508.06808v1#S8.SS2}.
\end{thebibliography}

\end{document}
